% !PRINTED 503
\weberpage{503}
%<<CONT
vanishes, but arbitrarily many, since one can multiply each
$\Phi(t)$ by an arbitrary function of the same form. It can
also happen that already a product of fewer than $n$ factors
of $N(t-\omega)$ acquires rational coefficients, from which
functions $\Phi(t)$ of degree lower than $n$ arise, which
vanish for $t = \omega$. From this it follows:

7. \gesperrt{For every functional $\omega$ of the field $\Omega$
there are infinitely many functions $\Phi(t)$ which
vanish for $t = \omega$, in which the highest power of $t$
has the coefficient 1, and whose remaining
coefficients are contained in $\overline{R}$.}

We then say that $\omega$ is a root of the equation
\[
\Phi(t) = 0.
\]

Among the various functions $\Phi(t)$ whose existence is
asserted in Theorem 7. there is one and only one $F(t)$ of
lowest possible degree. For if there existed two such
functions $F(t)$ and $F_1(t)$ of the same degree $m$, the
difference $F(t) - F_1(t)$ would be at most of degree
$m - 1$ with respect to $t$ and would vanish for $t = \omega$.
If one divides by the coefficient of the highest power of $t$,
one obtains a function $\Phi(t)$ of lower degree than $F(t)$,
which by the assumption about $F(t)$ does not exist.

The function $F(t)$ cannot be decomposed in the field
$\overline{R}$ into factors that are integral functions of $t$.
For if it split into several factors of the same form
$F_1(t)$, $F_2(t)$, then one of these factors, which are after
all of lower degree than $F(t)$ itself, would have to vanish
for $t = \omega$, contrary to our assumption.

Every function $\Phi(t)$ of Theorem 7. is then divisible by
this irreducible function $F(t)$, so that the quotient
$\Phi(t):F(t)$ is an integral function of $t$ in $\overline{R}$,
in which the highest power of $t$ has the coefficient 1
(cf. Bd. I, \S.\,141, I.).

Among the functions $\Phi(t)$ there is also, as already
remarked, the norm $N(t-\omega)$, and consequently
$N(t-\omega)$ is divisible by $F(t)$. If the two functions are
of the same degree, then $N(t-\omega) = F(t)$. But if the
degree of $F(t)$ is lower than $n$, the quotient
$N(t-\omega):F(t)$ vanishes at least for one more of the
functionals conjugate to $\omega$, and consequently for all of
them; consequently also for $t = \omega$, and hence this
quotient
%CONT>>
